% 原书物理页87--90，2.3.
\section{自同构与正规扩张}\label{sec:galois-2-3}
本节开始先固定 $k$ 的一个代数闭包 $\bar k$.
\subsection{$\bar k$ 的自同构}\label{sec:galois-2-3-1}
通常记
\begin{align*}
\operatorname{Aut}(\bar k/k):=\operatorname{Aut}_{k\text{-alg}}(\bar k)=\operatorname{Hom}_{k\text{-alg}}(\bar k,\bar k)
\end{align*}
为扩张 $\bar k\supset k$ 的自同构群. 若 $\bar k'$ 为 $k$ 的另一个代数闭包，则每个扩张同构 $\psi:\bar k'\xrightarrow{\sim}\bar k$ 都诱导同构
\begin{align*}
\sigma\longmapsto\psi^{-1}\sigma\psi:\operatorname{Aut}(\bar k/k)\xrightarrow{\sim}\operatorname{Aut}(\bar k'/k).
\end{align*}
\begin{lemma*}
设 $K\supset k$ 是 $k$ 的代数扩张，$\iota_1,\iota_2:K\hookrightarrow\bar k$ 为两个 $k$ 嵌入. 则存在自同构 $\sigma\in\operatorname{Aut}(\bar k/k)$，使得 $\iota_2=\sigma\circ\iota_1$.
\end{lemma*}
\begin{proof}
这是命题~\ref{sec:galois-2-2-3} 的推论. 事实上，嵌入 $\iota_2$ 给出 $K$ 的一个代数闭包，而嵌入 $\iota_1$ 使 $\bar k$ 成为 $K$ 的代数扩张. 命题~\ref{sec:galois-2-2-3} 于是给出 $K$ 扩张态射 $\sigma:\bar k\to\bar k$. 但须注意，左边经 $\iota_1$ 成为 $K$ 的扩张，右边则经 $\iota_2$. 因而按 $K$ 扩张态射的定义，$\sigma\circ\iota_1=\iota_2$. 又因为 $\iota_1$ 与 $\iota_2$ 均为 $k$ 线性的，$\sigma$ 也是如此. 故 $\sigma\in\operatorname{Aut}(\bar k/k)$.
\end{proof}
\subsection{$\bar k$ 中的共轭}\label{sec:galois-2-3-2}
\begin{proposition*}
对 $\alpha,\beta\in\bar k$，下列性质等价：
\begin{enumerate}[label=\roman*)]
\item 存在 $\sigma\in\operatorname{Aut}(\bar k/k)$，使 $\sigma(\alpha)=\beta$.
\item $f_\alpha=f_\beta$（在 $k$ 上的极小多项式）.
\end{enumerate}
这些性质成立时，称 $\alpha$ 与 $\beta$ \newterm{galois-conjugate-elements}{共轭}{conjugate elements}.
\end{proposition*}
\begin{proof}
对 $\sigma\in\operatorname{Aut}(\bar k/k)$，仍以 $\sigma:\bar k[X]\to\bar k[X]$ 表示延拓 $\sigma$ 且将 $X$ 映到 $X$ 的唯一 $k$ 代数自同构. 因而对任意 $f\in\bar k[X]$ 及 $x\in\bar k$，有 $\sigma(f(x))=\sigma(f)(\sigma(x))$. 又因为 $x$ 在 $k$ 上的极小多项式 $f_x$ 属于 $k[X]$，所以 $\sigma(f_x)=f_x$.

$i)\Rightarrow ii)$. 设 $\beta=\sigma(\alpha)$，其中 $\sigma\in\operatorname{Aut}(\bar k/k)$. 则
\begin{align*}
f_\alpha(\beta)=f_\alpha(\sigma(\alpha))=\sigma(f_\alpha)(\sigma(\alpha))=\sigma(f_\alpha(\alpha))=0.
\end{align*}
所以 $f_\beta\mid f_\alpha$. 同理 $f_\alpha\mid f_\beta$，最终 $f_\alpha=f_\beta$.

$ii)\Rightarrow i)$. 若 $f_\alpha=f_\beta$，则存在唯一的 $k$ 代数同构 $k[\alpha]\xrightarrow{\sim}k[\beta]$ 将 $\alpha$ 映到 $\beta$（通过 $k[X]/(f)$，其中 $f=f_\alpha=f_\beta$）. 再复合包含 $k[\beta]\subset\bar k$，得到将 $\alpha$ 映到 $\beta$ 的嵌入 $\iota:k[\alpha]\hookrightarrow\bar k$. 对 $K=k[\alpha]$、自然包含 $\iota_1$ 及 $\iota_2=\iota$ 应用前一引理，便得到延拓 $\iota$ 的嵌入 $\sigma:\bar k\to\bar k$，它将 $\alpha$ 映到 $\beta$.
\end{proof}
\begin{example*}
若 $k=\mathbb R$、$\bar k=\mathbb C$，则 $\operatorname{Aut}(\mathbb C/\mathbb R)=\{\mathrm{id},z\mapsto\bar z\}$，此命题中的共轭概念就成为通常的复共轭.
\end{example*}
\begin{example*}
设 $k=\mathbb Q$、$\bar k=\overline{\mathbb Q}$. $\sqrt[3]{2}$ 的共轭元素集合为 $\{\sqrt[3]{2},j\sqrt[3]{2},j^2\sqrt[3]{2}\}$，其中 $j=(-1+i\sqrt3)/2$ 是一个三次本原单位根.
\end{example*}
\subsection{$\bar k$ 的正规子扩张}\label{sec:galois-2-3-3}
\begin{proposition*}
设 $K$ 为 $\bar k$ 的子扩张. 下列性质等价：
\begin{enumerate}[label=\roman*)]
\item 对每个 $\alpha\in K$，$f_\alpha$ 在 $\bar k$ 中的所有根都属于 $K$.
\item 对每个 $\sigma\in\operatorname{Aut}(\bar k/k)$，有 $\sigma(K)\subset K$.
\end{enumerate}
若这些性质成立，就称 $K$ 为 $\bar k$ 的正规子扩张.
\end{proposition*}
\begin{proof}
$i)\Rightarrow ii)$. 设 $\sigma\in\operatorname{Aut}(\bar k/k)$、$\alpha\in K$. 前一命题说明 $\sigma(\alpha)$ 为 $f_\alpha$ 的根，因而假设 $i)$ 蕴涵 $\sigma(\alpha)\in K$，于是得到 $ii)$.

$ii)\Rightarrow i)$. 设 $\alpha\in K$，$\beta$ 为 $f_\alpha$ 的另一个根. 则 $f_\beta=f_\alpha$，由前一命题存在 $\sigma$ 使 $\sigma(\alpha)=\beta$. 由假设 $ii)$，$K$ 在 $\sigma$ 下稳定，故 $\beta\in K$.
\end{proof}
\begin{corollary*}
设 $K\subset\bar k$ 为正规子扩张. 限制映射 $\sigma\mapsto\sigma|_K$ 诱导满的群态射
\begin{align*}
\operatorname{Aut}(\bar k/k)\twoheadrightarrow\operatorname{Aut}(K/k),
\end{align*}
其核为 $\operatorname{Aut}(\bar k/K)$.
\end{corollary*}
\begin{proof}
前一命题的 $ii)$ 保证该映射有定义. 它显然是群态射. 最后，上一个引理保证 $K$ 的每个自同构 $\gamma$ 都能延拓为 $\bar k$ 的自同构 $\sigma$：只须取 $\iota_1$ 为自然包含，$\iota_2$ 为 $\gamma$ 与自然包含的复合，应用该引理即可. 因而此映射满. 其核由满足 $\sigma|_K=\mathrm{id}_K$ 的自同构 $\sigma\in\operatorname{Aut}(\bar k/k)$ 组成，亦即 $\bar k$ 的 $K$ 代数自同构，正如所述.
\end{proof}
\begin{example*}
设 $k=\mathbb Q$、$\bar k=\overline{\mathbb Q}$.
\begin{itemize}
\item $\mathbb Q[j]$ 是 $\mathbb Q$ 的正规扩张（次数为 $2$），因为每个 $\sigma\in\operatorname{Aut}(\overline{\mathbb Q})$ 都把 $j$ 映到 $j$ 或 $j^2$，因而使 $\mathbb Q[j]$ 稳定.
\item $\mathbb Q[\sqrt[3]{2}]$ 是 $\mathbb Q$ 的三次扩张，但不正规，因为存在 $\sigma\in\operatorname{Aut}(\overline{\mathbb Q})$，使 $\sigma(\sqrt[3]{2})=j\sqrt[3]{2}\notin\mathbb Q[\sqrt[3]{2}]$.
\item $\mathbb Q[j,\sqrt[3]{2}]$ 是 $\mathbb Q$ 的正规扩张（次数为 $6$）.
\end{itemize}
\end{example*}
\subsection{正规扩张}\label{sec:galois-2-3-4}
这里不再限于固定代数闭包 $\bar k$ 内的情形.
\begin{proposition*}
设 $K\supset k$ 是代数扩张. 下列性质等价：
\begin{enumerate}[label=\roman*)]
\item 对每个 $\alpha\in K$，多项式 $f_\alpha$ 在 $K[X]$ 中分裂.
\item 若 $\iota_1,\iota_2$ 为 $K$ 到一个代数闭包 $\bar k$ 的两个嵌入，则 $\iota_1(K)=\iota_2(K)$.
\item $K$ 在到任意代数闭包的任意嵌入下的像，都是前一分目意义下的正规子扩张.
\end{enumerate}
若这些性质成立，则称 $K\supset k$ 为\newterm{galois-normal-extension}{正规扩张}{normal extension}.
\end{proposition*}
\begin{proof}
$i)\Rightarrow ii)$. 设 $\alpha\in K$. 因为 $f_\alpha$ 分裂，每个嵌入 $\iota_1,\iota_2$ 都诱导 $f_\alpha$ 在 $K$ 中的根集到其在 $\bar k$ 中的根集的双射. 特别地，$\iota_2(\alpha)$ 为 $f_\alpha$ 的根，故形如 $\iota_1(\beta)$，从而属于 $\iota_1(K)$. 这证明 $\iota_2(K)\subset\iota_1(K)$；另一包含由对称性得到.

$ii)\Rightarrow iii)$. 设 $\iota:K\hookrightarrow\bar k$. 对任意 $\sigma\in\operatorname{Aut}(\bar k/k)$，由 $ii)$ 有 $\sigma(\iota(K))=\iota(K)$，故 $\iota(K)$ 为 $\bar k$ 的正规子扩张.

$iii)\Rightarrow i)$. 可将 $K$ 嵌入代数闭包，$\iota:K\hookrightarrow\bar k$. 此时多项式 $\iota(f_\alpha)$ 在 $\bar k[X]$ 中分裂，可写为 $\iota(f_\alpha)=\prod_{i=1}^m(X-x_i)^{v_i}$. 而 $iii)$ 蕴涵各 $x_i$ 属于 $\iota(K)$，设 $x_i=\iota(\alpha_i)$. 因而 $f_\alpha=\prod_i(X-\alpha_i)^{v_i}$ 在 $K[X]$ 中分裂.
\end{proof}
\begin{example*}
以``抽象''的方式（即不嵌入 $\overline{\mathbb Q}$）重述前一分目的例子：扩张 $\mathbb Q[X]/(X^2+X+1)\supset\mathbb Q$ 正规，而扩张 $\mathbb Q[X]/(X^3-2)\supset\mathbb Q$ 不正规.
\end{example*}
\subsection{多项式的分裂域}\label{sec:galois-2-3-5}
下面给出正规扩张的基本例子.
\begin{definition*}
设 $f\in k[X]$. $f$ 的一个\newterm{galois-splitting-field}{分裂域}{splitting field}是满足以下条件的 $k$ 扩张 $K$：
\begin{itemize}
\item $f$ 在 $K[X]$ 中分裂，即存在 $x_1,\ldots,x_n\in K$，使 $f=\lambda\prod_{i=1}^n(X-x_i)$；
\item $K$ 由 $f$ 的根 $x_i$ 生成.
\end{itemize}
\end{definition*}
\begin{corollary*}
每个多项式都有分裂域，且在（不唯一的）同构意义下唯一. 此域还是 $k$ 的正规扩张.
\end{corollary*}
\begin{proof}
设 $\bar k$ 为 $k$ 的一个代数闭包. 多项式 $f$ 分裂为 $f=\lambda\prod_{i=1}^n(X-x_i)$，其中 $\lambda\in k$，$x_1,\ldots,x_n\in\bar k$. 则 $\bar k$ 的子域 $K_f=k(x_1,\ldots,x_n)$ 是 $f$ 的分裂域. 再设 $K'\supset k$ 为 $f$ 的另一个分裂域. 则 $K'$ 在 $k$ 上代数，故可嵌入 $\bar k$. 其像由 $f$ 的根生成，因而等于 $K_f$. 所以 $K'$ 同构于 $K_f$. 这也说明 $K_f$ 正规.
\end{proof}
\begin{example*}
$X^3-2$ 在 $\mathbb Q$ 上的分裂域是 $\mathbb Q[j,\sqrt[3]{2}]$.
\end{example*}
\begin{exercise*}
设 $n,m\in\mathbb N$. 证明 $X^n-m$ 的分裂域为 $\mathbb Q[\zeta_n,\sqrt[n]{m}]$，其中 $\zeta_n=\exp(2i\pi/n)$，$\sqrt[n]{m}$ 为 $m$ 的唯一正实 $n$ 次方根.
\end{exercise*}
\noindent\textbf{译注.} 原文允许 $n,m\in\mathbb N$，但此式及``唯一正实根''要求 $n,m>0$；特别是 $m=0$ 时所述分裂域一般不成立. 上文分裂域定义及推论亦未明确排除零多项式.
\begin{remark*}
$\operatorname{Aut}(K_f/k)$ 在 $K_f$ 上的作用置换 $f$ 在 $K_f$ 中的根集 $f^{-1}(0)$. 因为这些根生成 $K_f$，故有到置换群的单射
\begin{align*}
\operatorname{Aut}(K_f/k)\hookrightarrow\mathfrak S_{f^{-1}(0)}.
\end{align*}
Galois 理论的基本想法是，用群 $\operatorname{Aut}(K_f/k)$ 作为代数方程 $f=0$ 的对称群. 然而这个群有时会是平凡群：取 $k=\mathbb F_p(T)$、$f=X^p-T$. 此时 $K_f=\mathbb F_p(T)[X]/(X^p-T)=\mathbb F_p(T^{1/p})$. 事实上，在 $K_f$ 中有分解 $X^p-T=(X-T^{1/p})^p$，说明 $T^{1/p}$ 是 $T$ 唯一的 $p$ 次方根（重数为 $p$）. 所以群 $\mathfrak S_{f^{-1}(0)}$ 平凡，$\operatorname{Aut}(K_f/k)$ 亦然. 这一称为``不可分性''的现象将在后面几节研究.
\end{remark*}



